\chapter{Tetany}

\section*{Definition}
Tetany is a neurological syndrome of neuromuscular hyperexcitability characterized by distal paresthesias, muscle cramps, carpopedal spasm, and laryngospasm, most commonly resulting from hypocalcemia. Latent tetany may be elicited by provocative maneuvers such as Chvostek or Trousseau signs.

\section*{What Happens in Your Body}
Extracellular ionized calcium stabilizes the voltage-gated sodium channel in its resting (inactivated) state, raising the threshold for action potential generation. Hypocalcemia reduces the membrane-stabilizing effect, allowing spontaneous depolarization of peripheral nerves. This increases the frequency of repetitive firing in motor and sensory axons, producing paresthesias and involuntary muscle contractions. Hypomagnesemia impairs PTH secretion and end-organ responsiveness, exacerbating hypocalcemia. Alkalosis (respiratory or metabolic) increases calcium binding to albumin, lowering the ionized fraction and producing tetany even with normal total calcium.

\section*{Causes}
\begin{itemize}
  \item \textbf{Hypocalcemia:} Renal failure, hypoparathyroidism (post-surgical, autoimmune, DiGeorge syndrome), vitamin D deficiency, pancreatitis, rhabdomyolysis, osteoporosis medications (denosumab, bisphosphonates), hungry bone syndrome after parathyroidectomy, tumor lysis syndrome
  \item \textbf{Hypomagnesemia:} Chronic alcoholism, diuretic use, proton pump inhibitor use, Gitelman syndrome, cisplatin therapy, aminoglycosides
  \item \textbf{Alkalosis:} Hyperventilation (anxiety, panic attack), vomiting, primary hyperaldosteronism, Cushing syndrome, salicylate poisoning early phase
  \item \textbf{Hypokalemia:} Can potentiate tetany though less common as direct cause
  \item \textbf{Other:} Hypocalcemic tetany of the newborn, hypoproteinemia, critical illness
\end{itemize}

\section*{When to See a Doctor}
See your doctor if:
\begin{itemize}
  \item Carpopedal spasm (flexion of wrist, MCP, and IP joints with thumb adduction)
  \item Perioral or digital paresthesias with muscle cramping
  \item Laryngospasm causing high-pitched breathing sound or difficulty breathing
  \item Seizure, altered mental status, or prolonged QT interval on ECG
\end{itemize}

\section*{Related Symptoms}
\begin{itemize}
  \item Perioral and distal fingertip paresthesias (earliest symptom)
  \item Carpopedal spasm: wrist flexion, thumb adduction, finger extension at IP joints
  \item Laryngospasm: high-pitched breathing sound, hoarseness, respiratory distress
  \item Chvostek sign: facial nerve tap triggers ipsilateral facial twitching
  \item Trousseau sign: BP cuff inflation to above systolic for 3 minutes induces carpal spasm
  \item Generalized muscle cramps, stiffness, and fatigue
\end{itemize}
\textbf{Important:} Chvostek sign is present in 10--25\% of normal individuals; Trousseau sign is more specific (sensitivity 94\%) for latent tetany and should be preferentially used.

\section*{Self-Care / Home Management}
\begin{itemize}
  \item Tetany is potentially serious; evaluation and treatment of the underlying cause is required
  \item For mild, recurrent tetany from chronic hypoparathyroidism: oral calcium carbonate (1--2 g/day) and calcitriol (0.25--0.5 mcg/day)
  \item Stress reduction and breathing techniques for hyperventilation-induced tetany
  \item Avoidance of alcohol and PPI medications if hypomagnesemic
  \item Maintain adequate dietary calcium, magnesium, and vitamin D
\end{itemize}

\section*{How Your Doctor Will Evaluate You}
\begin{itemize}
  \item Serum calcium (total and ionized) as primary screening
  \item Serum magnesium, phosphate, potassium, and albumin for context
  \item PTH level (low in hypoparathyroidism, high in vitamin D deficiency/renal failure)
  \item Vitamin D (25-OH and 1,25-OH) levels
  \item Renal function (BUN, creatinine) and pancreatic enzymes if pancreatitis suspected
  \item ECG for QT interval prolongation (corrected QT > 440 ms suggests severe hypocalcemia)
  \item Arterial blood gas if alkalosis-driven tetany is suspected
\end{itemize}

\section*{Treatment Options}
\begin{itemize}
  \item Acute severe tetany: IV calcium gluconate 10 mL of 10\% solution over 10--20 minutes with cardiac monitoring
  \item Maintenance: Oral calcium carbonate 1--2 g/day in divided doses
  \item Hypomagnesemia: IV or oral magnesium sulfate (2--4 g IV over 20 minutes for severe)
  \item Vitamin D supplementation: Calcitriol 0.25--0.5 mcg/day for hypoparathyroidism
  \item Treatment of underlying cause: pancreatic enzymes, hemodialysis if renal failure, stopping offending medication
  \item Hyperventilation tetany: Rebreathing into a paper bag (increases CO$_2$, reduces alkalosis), benzodiazepines for severe panic-driven hyperventilation
  \item Long-term monitoring and follow-up with endocrinology for chronic hypoparathyroidism
\end{itemize}

\clearpage
